\section{Introduction}

Machine-learning asset-pricing models can use many firm characteristics while producing a small number of latent factors or concentrated hidden representations. This raises an identification question: does a compact fitted representation reveal an economic dimension, or does it summarize the information that matters for a particular sample and evaluation criterion? The distinction matters when a selected network width is interpreted as the number of priced risks or used to fix model complexity over time.

Flexible return predictors and characteristic-based factor models provide several ways to summarize the cross-section \citep{gu2020empirical,kelly2019characteristics,gu2021autoencoder,chen2024deep}. Their statistical and economic objects differ. A low effective rank of activations measures concentration; a small factor set concerns a payoff span; a minimum validation loss identifies a preferred candidate under a specified criterion. None of these quantities is automatically the smallest number of structural shocks. We examine the additional evidence needed to connect them. We do not attribute a common structural claim to all of the cited studies.

Our design separates four questions: whether representations concentrate, whether the smallest pricing-sufficient dimension is identified, whether the conclusion persists across measurements and periods, and whether a historically feasible decision rule produces economic gains. These are distinct tests rather than interchangeable measures of model success. The pricing analysis uses six nominal dimensions, five fixed seeds, two archived characteristic pipelines, and several test-asset families. A known-loading simulation isolates statistical selection error. A separate return-prediction replay tests whether retrospective changes in updating value are usable in real time.

The results show substantial dependence on the evaluation design. In development, the ranking of six dimensions changes across weighting geometries and characteristic pipelines; the average rank correlations between the pipelines are negative. Their comparison is descriptive: it combines construction and coverage differences with a one-month difference in return-window timing. Core-86's locked 2020--2025 evaluation retains only one of four prespecified directional criteria, and its market-attenuation intervals include zero. The networks were unchanged, but legacy development had included January 2020, so this later evaluation is not an entirely untouched holdout for the research process.

The simulation provides a statistical explanation for one possible source of instability. Even with known loadings and covariance, an independent noisy evaluation mean can favor an unnecessary direction. We derive this result for a finite nested Gaussian family when training and evaluation samples grow proportionately. A suitable vanishing penalty restores consistency in that design. The argument illustrates a familiar distinction between predictive accuracy and exact model selection \citep{shao1993cv,yang2007cv}; it is not a general theorem for trained neural networks. In the simulation grid, average exact recovery is 64.2\% under population-mean evaluation and 32.6\% under noisy evaluation.

Representation concentration nevertheless appears throughout the frozen network sample. Median participation rank declines from 20.25 to 9.05 across the three hidden layers. Yet final-layer rank adds little to held-out prediction of model-level pricing loss after nominal factor count is controlled for. The separate temporal exercise also distinguishes retrospective patterns from usable rules: updating becomes more aligned with subsequent errors in parts of the later sample, while the tested detector reacts late and the six primary incremental portfolio comparisons fail multiplicity-adjusted significance tests. These return predictors are independent of the pricing-dimension networks; they provide supporting evidence about updating, not a direct test of time variation in the latent pricing dimension.

The contribution is an empirical comparison of these measurement and decision criteria within one documented research sequence. Test assets and weighting matter for pricing conclusions \citep{hansen1997assessing,lewellen2010skeptical}; factor searches require attention to replication and multiplicity \citep{feng2020taming,hou2020replicating}; model comparisons can change after transaction costs \citep{detzel2023costs}. Our evidence connects these issues to the interpretation of compact neural representations. It supports reporting the range of dimensions compatible with a stated loss tolerance, together with the measurement, chronology, and economic-use limitations of that range. It does not establish that a particular alternative selector or monitoring system performs better.

Section~2 develops the framework and the known-loading result. Section~3 describes data, estimators, chronology, and inference. Section~4 reports dimension sensitivity and the simulation. Sections~5 and~6 examine updating and representation geometry, respectively. Section~7 discusses the scope and limitations of the findings.
